\subsection{Coin Task}

For the coin task, we consider sequences of the form
\[
x^{(k,b)}_{1:t}
= x^{(k,b)}_1,
x^{(k,b)}_2,
\ldots,
x^{(k,b)}_t,
\]
where $k$ indexes the task and $b$ indexes the batch.
From layer $l$ we extract the hidden state at the $t$-th token and denote
it by $\vh^{(l,b)}_{k,t}$.
Suppressing the layer index, we adopt the representation model
\[
\vh^{(b)}_{k,t}
= \vmu_t
+ \vtheta_{k,t}\!\left(x^{(k,b)}_t\right)
+ \vepsilon^{(b)}_{k,t}\!\left(x^{(k,b)}_{1:t}\right).
\]
The current token $x^{(k,b)}_t$ may now depend on the
task identity $k$.
Consequently, task information and token information are generally entangled at the level
of the hidden representation, and the task index cannot be removed from
$\vtheta_{k,t}(\cdot)$.

Our analysis focuses on the two properties of this decomposition. 
\paragraph{(P1) Vanishing residual variance.}
For the hidden representation to be well approximated by the task--token component
$\vtheta_{k,t}(x^{(k,b)}_t)$, the residual term must behave as noise across batches.
Specifically, we require
\[
\operatorname{Var}_b\!\left[\vepsilon^{(b)}_{k,t}(x_{1:t}^{(k,b)})\right] \to 0,
\]
as $t \to \infty$.
This reflects the fact that, in the coin flipping task, all task-relevant information needed for
next-token prediction is contained in the current covariate $x_t$, so dependence on earlier
context becomes asymptotically negligible.

To empirically assess (P1), we group samples across batches that share the same current token
value $x_t$ and form the conditional average
\[
\bar{\vh}_{k,t}(x_t)
:= \frac{1}{|\mathcal{B}(x_t)|}
\sum_{b \in \mathcal{B}(x_t)} \vh^{(b)}_{k,t},
\]
where $\mathcal{B}(x_t) = \{ b : x^{(b)}_t = x_t \}$.
We then consider the residual
\[
\vh^{(b)}_{k,t} - \bar{\vh}_{k,t}(x_t),
\]
whose empirical variance estimates the conditional variance of
$\vepsilon^{(b)}_{k,t}$ given $(k, x_t)$.



\paragraph{(P2) Linear task--token decomposition.}
Based on the belief state hypothesis that $\vh_{k,t}^{(b)}$ should encode the current token information and the posterior
$$
\alpha^{(b)}_{t,k} = \P\left(Z = k \mid x^{(b)}_{1:t}\right),
$$
we would expect that there exists a map $g$ such that 
\begin{equation*}
    g(\vtheta_{k,t}(x)) = \left(\alpha^{(b)}_{t,1:K}, x\right). 
\end{equation*}
When $g$ is a linear map and nonlinear interactions between task identity and the current token are weak,
we would further posit the approximation
\[
\vtheta_{k,t}(x) \approx \vnu_{k,t} + \vphi_t(x),
\]
where $\vnu_{k,t}$ is a task-specific vector that encodes the belief states $\alpha^{(b)}_{t}$ and $\vphi_t(x)$ depends only on the token.
We refer to this as a \emph{linear task--token decomposition}.
Failure of this approximation indicates the presence of genuine token-dependent task representations.




To empirically assess property (P2), let
\[
\bar{\vh}_{k,t} := \frac{1}{B}\sum_{b=1}^B \vh^{(b)}_{k,t}
\]
denote the task-averaged hidden state.
%Then
%\[
%\vh^{(b)}_{k,t} - \bar{\vh}_{k,t}
%\;\approx\;
%\vtheta_{k,t}\!\left(x^{(b)}_t\right)
%- \frac{1}{B}\sum_{b=1}^B \vtheta_{k,t}\!\left(x^{(b)}_t\right) \approx \vphi_t(x_t^{(b)}) - \frac{1}{B} \sum_{b=1}^B \vphi_t(x_t^{(b)}),
%\]
%whose variance generally does not vanish, since the covariates $x^{(b)}_t$ may differ
%across batches.
To isolate task-level structure, we instead consider the task-centered quantity
\[
\begin{aligned}
\vh^{(b)}_{k,t} - \frac{1}{K}\sum_{k=1}^K \vh^{(b)}_{k,t}
&\approx
\vtheta_{k,t}\!\left(x^{(b)}_t\right)
- \frac{1}{K}\sum_{k=1}^K \vtheta_{k,t}\!\left(x^{(b)}_t\right)\\
&\approx \vnu_{k,t} - \frac{1}{K} \sum_{k=1}^K \vnu_{k,t}. 
\end{aligned}
\]
Under a linear task--token decomposition, the token-dependent terms cancel after averaging,
and the variance of this quantity vanishes asymptotically as $t \to \infty$.
Persistent variance therefore provides direct evidence of nonlinear task--token interactions.